\begin{abstract}
2026年云栖大会期间，浙江大学联合阿里云开展 Qiushi Engine（求是引擎）
72小时自主科研直播，公开探索高维接吻数问题。
这一问题研究彼此不重叠的等球同时接触中心球的最大数量。
扩大已有构型，需要在严密的距离约束中发现新的几何可能。
Qiushi Engine 通过持续自主探索，取得\ResultCount{}个维数的接吻数改进下界。
结果覆盖32至39维、43维和45维。

这些突破体现了一个共同的数学思想：大型构型的增长空间，存在于各部分的组合方式、
它们在母结构中的位置，以及能够揭示其规模的数学关系之中。
Qiushi Engine 将编码理论、格几何与组合结构联系起来，通过联合替换扩大局部构型，
从另一格壳中找到与既有点集相容的新方向，并利用矩恒等式从庞大母壳中确定或估计
所需截面的规模。其中，45维的构造以298个显式见证推出7380090点下界，
将数千万个向量的计数问题化为简短的精确论证。

这些发现展示了 Qiushi Engine 围绕开放目标持续开展研究、建立跨领域的数学联系，
并将新构想发展为构造与证明的能力。本次挑战以新的数学成果呈现了持续自主探索
的创造力。公开直播让这一过程走近公众，开放的构造、论证与精确数据则让科学
共同体能够重建并继续发展这些发现。

\par\smallskip\noindent\textbf{关键词}\quad
自主数学发现；接吻数；常重码；Hadamard变换；Leech格；格截面；球面设计
\end{abstract}
